\documentclass[11pt]{report}

\usepackage[T1]{fontenc}
\usepackage[utf8]{inputenc}
\usepackage{lmodern}
\usepackage[a4paper,margin=1in]{geometry}
\usepackage{amsmath,amssymb,amsthm,amscd,mathtools,mathrsfs}
\usepackage{enumitem}
\usepackage{tikz-cd}
\usepackage{xcolor}
\usepackage{microtype}
\usepackage{xurl}
\usepackage{hyperref}
\usepackage[nameinlink,capitalise,noabbrev]{cleveref}

\hypersetup{
  colorlinks=true,
  linkcolor=blue!50!black,
  citecolor=green!40!black,
  urlcolor=blue!60!black
}

% Visible dependency metadata.  Every local dependency is printed together
% with its numbered cross-reference, so the proof graph remains inspectable
% in both the TeX source and the compiled PDF.
\newcommand{\lean}[1]{%
  \par\smallskip\noindent\textcolor{blue!55!black}{\textsf{Lean:} }%
  \footnotesize\sloppy\nolinkurl{#1}\par}
\ExplSyntaxOn
\bool_new:N \l__weil_first_dependency_bool
\cs_new_protected:Npn \weil_audit_dependency_list:n #1
  {
    \tl_if_blank:nTF {#1}
      { \textit{none~(foundational~declaration)} }
      {
        \bool_set_true:N \l__weil_first_dependency_bool
        \clist_map_inline:nn {#1}
          {
            \bool_if:NTF \l__weil_first_dependency_bool
              { \bool_set_false:N \l__weil_first_dependency_bool }
              { ,\penalty0\space }
            \nolinkurl{##1}~[\ref{##1}]
          }
      }
  }
\NewDocumentCommand{\auditdependencylist}{m}
  { \weil_audit_dependency_list:n {#1} }
\ExplSyntaxOff
\newcommand{\uses}[1]{%
  \par\smallskip\noindent\textcolor{purple!60!black}{\textsf{Uses:} }%
  \footnotesize\sloppy\auditdependencylist{#1}\par}
\newcommand{\leanok}{\par\smallskip\noindent
  \textcolor{green!45!black}{\textsf{Lean status: theorem intended for formalization.}}\par}
\newcommand{\mathlibok}{\par\smallskip\noindent
  \textcolor{green!45!black}{\textsf{Library status: verified only after the pinned-contract audit.}}\par}
\newcommand{\formalizationnote}[1]{%
  \begin{quote}\small\raggedright\textbf{Formalization note.} #1\end{quote}}
\newcommand{\statusnote}[1]{%
  \begin{quote}\small\raggedright\textbf{Status.} #1\end{quote}}
\newcommand{\constructedby}[1]{%
  \par\smallskip\noindent\textcolor{purple!60!black}{\textsf{Constructed by:} }%
  \footnotesize\sloppy\auditdependencylist{#1}\par}
\newcommand{\openobligation}[1]{%
  \par\smallskip\noindent\textcolor{red!65!black}{\textsf{Open obligation:} }%
  \footnotesize\sloppy\nolinkurl{#1}\par}

\theoremstyle{definition}
\newtheorem{definition}{Definition}[section]
\newtheorem{construction}[definition]{Construction}
\newtheorem{example}[definition]{Example}

\theoremstyle{plain}
\newtheorem{theorem}[definition]{Theorem}
\newtheorem{proposition}[definition]{Proposition}
\newtheorem{lemma}[definition]{Lemma}
\newtheorem{corollary}[definition]{Corollary}

\theoremstyle{remark}
\newtheorem{remark}[definition]{Remark}
\newtheorem{warning}[definition]{Warning}
\newcommand{\Qell}{\mathbf{Q}_{\ell}}
\newcommand{\Zell}{\mathbf{Z}_{\ell}}
\newcommand{\Fq}{\mathbf{F}_{q}}
\newcommand{\Frob}{\operatorname{Frob}}
\newcommand{\Tr}{\operatorname{Tr}}
\newcommand{\wt}{\operatorname{wt}}
\newcommand{\coeff}{\operatorname{coeff}}

\begin{document}
\sloppy
\setcounter{chapter}{2}

\chapter{Zeta functions and finite Frobenius point counting}
\label{chap:weil-zeta-point-counting}

Unless a later scope paragraph explicitly strengthens the hypotheses, let $R$
be an arbitrary commutative ring.  We write $R[[T]]$ for the formal
power-series ring, $[T^n]A$ for the coefficient of $T^n$ in $A$, and
$c(A)=[T^0]A$ for its constant coefficient.  No topology or analytic
convergence is used: every construction below is coefficientwise finite.
Unless explicitly stated otherwise, all coefficient, degree, truncation,
period, divisor, and iteration indices in this chapter lie in $\mathbf N$.
The coefficient-ring hypothesis of a declaration is the most recent explicit
scope paragraph above it; each reset below is therefore part of the formal
statement envelope and is not merely expository prose.
All Lean declarations in the first four sections, from formal power-series
calculus through finite Frobenius orbits, that use
\texttt{PowerSeries.subst}, \texttt{PowerSeries.derivative},
\texttt{Set.Finite.toFinset}, \texttt{Set.ncard}, or
\texttt{Quotient.out} are made in a \texttt{noncomputable section}.
Proofs and finset constructions that require decidable equality begin with
\texttt{classical}; no finiteness or decidability instance is being suppressed
as mathematical input.

\section{Finite geometric-Frobenius orbits}
\label{sec:weil-finite-orbits}

No coefficient ring is in scope in the combinatorial declarations of this
section.  When Euler products resume later, their rational coefficient ring
is stated there explicitly.

The letter $F$ in this section denotes a permutation modeling
\emph{geometric} Frobenius; no arithmetic Frobenius occurs here.  To keep
natural iteration distinct from powers in the permutation group, write
$F^{[n]}(x)$ for
$\texttt{Function.iterate }(F:\Omega\to\Omega)\,n\,x$ when
$n\in\mathbf N$.  The notation $F^z$, used only for $z\in\mathbf Z$, denotes
the corresponding group power in $\operatorname{Equiv.Perm}(\Omega)$.

\begin{definition}[Finite-orbit geometric-Frobenius datum]
  \label{def:weil-finite-orbit-frobenius-datum}
  \uses{}
  \lean{Weil.FiniteOrbitFrobeniusData}

  A finite-orbit geometric-Frobenius datum $D$ is a structure with fields:
  \begin{enumerate}[label=\textup{(\roman*)}]
    \item \texttt{points : Type u} and
      \texttt{frobenius : Equiv.Perm points}, written $\Omega$ and $F$;
    \item \texttt{period : $\Omega\to\mathbf N$},
      \texttt{period\_pos : $\forall x,0<\texttt{period}(x)$}, and
      \texttt{iterate\_period : $\forall x,
      F^{[\texttt{period}(x)]}(x)=x$};
    \item \texttt{positiveFixedPoints :
      $\forall n:\mathbf N,0<n\to\operatorname{Finset}\Omega$} and
      \texttt{mem\_positiveFixedPoints}, asserting for $h:0<n$ and every
      $x$ the equivalence
      \[
        x\in\texttt{positiveFixedPoints }n\,h
          \quad\Longleftrightarrow\quad F^{[n]}(x)=x.
      \]
  \end{enumerate}
\end{definition}

\paragraph{Formal datum context.}
For every declaration from here to the end of this section, fix an arbitrary
$D:\texttt{Weil.FiniteOrbitFrobeniusData}$ and write
$\Omega:=D.\texttt{points}$ and $F:=D.\texttt{frobenius}$.  Every such
declaration is universally quantified over this explicit parameter $D$.

\begin{definition}[Total fixed-point finset]
  \label{def:weil-frobenius-fixed-point-finset}
  \uses{def:weil-finite-orbit-frobenius-datum}
  \lean{Weil.FiniteOrbitFrobeniusData.fixedPointFinset}

  For every $n\in\mathbf N$, define
  \[
    \operatorname{Fix}_D(n):=
    \begin{cases}
      \texttt{positiveFixedPoints }n\,h,&h:0<n,\\
      \varnothing,&n=0.
    \end{cases}
  \]
  Proof irrelevance makes the positive branch independent of the chosen
  witness $h$.
\end{definition}

\begin{lemma}[Membership in the total fixed-point finset]
  \label{lemma:weil-frobenius-fixed-point-finset-mem}
  \uses{def:weil-finite-orbit-frobenius-datum,
    def:weil-frobenius-fixed-point-finset}
  \lean{Weil.FiniteOrbitFrobeniusData.mem_fixedPointFinset}

  If $0<n$, then for every $x\in\Omega$,
  \[
    x\in\operatorname{Fix}_D(n)\quad\Longleftrightarrow\quad
    F^{[n]}(x)=x.
  \]
\end{lemma}

\begin{proof}
  \uses{def:weil-finite-orbit-frobenius-datum,
    def:weil-frobenius-fixed-point-finset}
  Unfold Definition~\ref{def:weil-frobenius-fixed-point-finset}.  The
  positive branch is \texttt{positiveFixedPoints n h}; apply the datum field
  \texttt{mem\_positiveFixedPoints}.  Proof irrelevance identifies its
  stored positivity witness with the witness used to unfold the definition.
\end{proof}

\begin{definition}[Frobenius orbit length]
  \label{def:weil-frobenius-orbit-length}
  \uses{def:weil-finite-orbit-frobenius-datum}
  \lean{Weil.FiniteOrbitFrobeniusData.orbitLength}

  For $x\in\Omega$, define
  \[
    d_F(x):=\texttt{Function.minimalPeriod }F\,x.
  \]
\end{definition}

\begin{lemma}[Positivity of Frobenius orbit length]
  \label{lemma:weil-frobenius-orbit-length-positive}
  \uses{def:weil-finite-orbit-frobenius-datum,
    def:weil-frobenius-orbit-length}
  \lean{Weil.FiniteOrbitFrobeniusData.orbitLength_pos}

  For every $x\in\Omega$, one has $d_F(x)>0$.
\end{lemma}

\begin{proof}
  \uses{def:weil-finite-orbit-frobenius-datum,
    def:weil-frobenius-orbit-length}
  The fields
  $(\texttt{period}(x),F^{[\texttt{period}(x)]}(x)=x)$ stored in the
  Frobenius datum, together with \texttt{period\_pos}, prove that
  $x\in\texttt{Function.periodicPts }F$.  Apply
  \texttt{Function.minimalPeriod\_pos\_of\_mem\_periodicPts} with
  $f=F$ and the point $x$.
\end{proof}

\begin{lemma}[Fixed-point criterion from the minimal period]
  \label{lemma:weil-frobenius-fixed-iff-orbit-length-divides}
  \uses{def:weil-frobenius-orbit-length}
  \lean{Weil.FiniteOrbitFrobeniusData.isPeriodicPt_iff_dvd}

  For $n\geq1$ and $x\in\Omega$,
  \[
    F^{[n]}(x)=x\quad\Longleftrightarrow\quad d_F(x)\mid n.
  \]
\end{lemma}

\begin{proof}
  \uses{def:weil-frobenius-orbit-length}
  The left side is
  \nolinkurl{Function.IsPeriodicPt F n x}.  Apply the exact mathlib theorem
  \nolinkurl{Function.isPeriodicPt_iff_minimalPeriod_dvd} with
  $f=F$, period $n$, and point $x$.  Replace
  \nolinkurl{Function.minimalPeriod F x} by $d_F(x)$ using the definition.
\end{proof}

\begin{definition}[Cyclic group generated by Frobenius]
  \label{def:weil-frobenius-cyclic-subgroup}
  \uses{def:weil-finite-orbit-frobenius-datum}
  \lean{Weil.FiniteOrbitFrobeniusData.cyclicSubgroup}

  Regard $F$ as an element of the group
  $\operatorname{Equiv.Perm}(\Omega)$.  Define
  \[
    \Gamma_F:=\texttt{Subgroup.zpowers }F.
  \]
  It acts on $\Omega$ by evaluation of permutations.
\end{definition}

\begin{definition}[Type of Frobenius orbits]
  \label{def:weil-frobenius-orbit-type}
  \uses{def:weil-frobenius-cyclic-subgroup}
  \lean{Weil.FiniteOrbitFrobeniusData.Orbit}

  Define the type of Frobenius orbits by
  \[
    \operatorname{Orb}_F(\Omega)
      :=\texttt{MulAction.orbitRel.Quotient }\Gamma_F\,\Omega.
  \]
  For $x\in\Omega$, write $[x]_F$ for its quotient class and, for an orbit
  $O$, write $|O|\subseteq\Omega$ for
  \nolinkurl{MulAction.orbitRel.Quotient.orbit O}.
\end{definition}

\begin{definition}[Selected representative of a Frobenius orbit]
  \label{def:weil-frobenius-orbit-representative}
  \uses{def:weil-frobenius-orbit-type}
  \lean{Weil.FiniteOrbitFrobeniusData.Orbit.representative}

  For $O\in\operatorname{Orb}_F(\Omega)$, define
  \[
    \operatorname{rep}(O):=\texttt{Quotient.out }O\in\Omega.
  \]
\end{definition}

\begin{lemma}[The selected representative represents its orbit]
  \label{lemma:weil-frobenius-orbit-representative-class}
  \uses{def:weil-frobenius-orbit-type,
    def:weil-frobenius-orbit-representative}
  \lean{Weil.FiniteOrbitFrobeniusData.Orbit.mk_representative}

  For every orbit $O$, one has
  $[\operatorname{rep}(O)]_F=O$ and
  $\operatorname{rep}(O)\in|O|$.
\end{lemma}

\begin{proof}
  \uses{def:weil-frobenius-orbit-representative}
  The first equality is the exact theorem \texttt{Quotient.out\_eq'} for the
  setoid \texttt{MulAction.orbitRel $\Gamma_F$ $\Omega$}.  Substitute this
  equality into
  \texttt{MulAction.orbitRel.Quotient.mem\_orbit}; that theorem states that
  $x\in|O|$ if and only if $[x]_F=O$, and gives the second assertion.
\end{proof}

\begin{definition}[Intrinsic degree of a Frobenius orbit]
  \label{def:weil-frobenius-orbit-degree}
  \uses{def:weil-frobenius-orbit-type}
  \lean{Weil.FiniteOrbitFrobeniusData.Orbit.degree}

  For $O\in\operatorname{Orb}_F(\Omega)$, define
  \[
    \deg_F(O):=\operatorname{ncard}(|O|).
  \]
  The finiteness and positivity of $|O|$ are proved below rather than included
  silently in this definition.
\end{definition}

\begin{lemma}[A Frobenius orbit is finite]
  \label{lemma:weil-frobenius-orbit-finite}
  \uses{def:weil-finite-orbit-frobenius-datum,
    def:weil-frobenius-cyclic-subgroup,
    def:weil-frobenius-orbit-type}
  \lean{Weil.FiniteOrbitFrobeniusData.Orbit.finite}

  For every $O\in\operatorname{Orb}_F(\Omega)$, the set $|O|$ is finite.
\end{lemma}

\begin{proof}
  \uses{def:weil-frobenius-orbit-representative,
    lemma:weil-frobenius-orbit-length-positive,
    lemma:weil-frobenius-orbit-representative-class}
  Put $x=\operatorname{rep}(O)$ and $d=d_F(x)>0$.  The exact equivalence
  \nolinkurl{MulAction.orbitZPowersEquiv F x} identifies
  \nolinkurl{MulAction.orbit (Subgroup.zpowers F) x} with
  $\mathbf Z/d\mathbf Z=\texttt{ZMod d}$.  Positivity of $d$ supplies the
  \texttt{NeZero d} instance, so \texttt{ZMod d} is finite.  Transport
  finiteness across the equivalence.  Finally,
  \nolinkurl{MulAction.orbitRel.Quotient.orbit_eq_orbit_out}, with
  \nolinkurl{Quotient.out_eq'}, identifies that orbit with $|O|$.
\end{proof}

\begin{lemma}[Orbit degree equals the minimal period of a representative]
  \label{lemma:weil-frobenius-orbit-degree-representative}
  \uses{def:weil-frobenius-orbit-degree,
    def:weil-frobenius-orbit-representative,
    def:weil-frobenius-orbit-length}
  \lean{Weil.FiniteOrbitFrobeniusData.Orbit.degree_eq_minimalPeriod}

  For every orbit $O$,
  \[
    \deg_F(O)=d_F(\operatorname{rep}(O)).
  \]
\end{lemma}

\begin{proof}
  \uses{lemma:weil-frobenius-orbit-finite,
    lemma:weil-frobenius-orbit-representative-class}
  Install the finite type structure on $|O|$ obtained from
  Lemma~\ref{lemma:weil-frobenius-orbit-finite}.  Rewrite $|O|$ as the orbit
  of $x=\operatorname{rep}(O)$ using
  \texttt{MulAction.orbitRel.Quotient.orbit\_eq\_orbit\_out}.  Apply
  \texttt{MulAction.minimalPeriod\_eq\_card} with the group element
  $a=F$ and point $b=x$.  It gives
  \[
    d_F(x)=\operatorname{card}
       (\texttt{MulAction.orbit (Subgroup.zpowers F) x}).
  \]
  For a finite set, \texttt{Set.ncard} is that subtype cardinal.  Reverse the
  displayed equality to match the definition of $\deg_F(O)$.
\end{proof}

\begin{lemma}[Orbit degree is positive]
  \label{lemma:weil-frobenius-orbit-degree-positive}
  \uses{def:weil-frobenius-orbit-degree}
  \lean{Weil.FiniteOrbitFrobeniusData.Orbit.degree_pos}

  For every Frobenius orbit $O$, one has $\deg_F(O)>0$.
\end{lemma}

\begin{proof}
  \uses{lemma:weil-frobenius-orbit-degree-representative,
    lemma:weil-frobenius-orbit-length-positive}
  Replace $\deg_F(O)$ by the minimal period of its representative and apply
  Lemma~\ref{lemma:weil-frobenius-orbit-length-positive}.
\end{proof}

\begin{lemma}[Degree of the orbit containing a point]
  \label{lemma:weil-frobenius-orbit-degree-point}
  \uses{def:weil-frobenius-orbit-type,
    def:weil-frobenius-orbit-degree,
    def:weil-frobenius-orbit-length}
  \lean{Weil.FiniteOrbitFrobeniusData.degree_mk}

  For every $x\in\Omega$,
  \[
    \deg_F([x]_F)=d_F(x).
  \]
\end{lemma}

\begin{proof}
  \uses{lemma:weil-frobenius-orbit-finite,
    lemma:weil-frobenius-orbit-degree-representative}
  The quotient-orbit theorem
  \nolinkurl{MulAction.orbitRel.Quotient.orbit_mk} identifies
  $|[x]_F|$ with
  \nolinkurl{MulAction.orbit (Subgroup.zpowers F) x}.  Install its finite
  type structure using Lemma~\ref{lemma:weil-frobenius-orbit-finite} and apply
  \nolinkurl{MulAction.minimalPeriod_eq_card} with $a=F,b=x$.  Converting
  the cardinal to \texttt{Set.ncard} gives the displayed equality directly.
\end{proof}

\begin{lemma}[Fixed-point criterion in terms of intrinsic orbit degree]
  \label{lemma:weil-frobenius-fixed-iff-orbit-degree-divides}
  \uses{def:weil-frobenius-orbit-type,
    def:weil-frobenius-orbit-degree}
  \lean{Weil.FiniteOrbitFrobeniusData.fixed_iff_degree_dvd}

  For $n\geq1$ and $x\in\Omega$,
  \[
    F^{[n]}(x)=x\quad\Longleftrightarrow\quad
    \deg_F([x]_F)\mid n.
  \]
\end{lemma}

\begin{proof}
  \uses{lemma:weil-frobenius-fixed-iff-orbit-length-divides,
    lemma:weil-frobenius-orbit-degree-point}
  Replace $\deg_F([x]_F)$ with $d_F(x)$ by
  Lemma~\ref{lemma:weil-frobenius-orbit-degree-point}.  The result is exactly
  Lemma~\ref{lemma:weil-frobenius-fixed-iff-orbit-length-divides}.
\end{proof}
\begin{definition}[Explicit finset of points in an orbit]
  \label{def:weil-frobenius-orbit-point-finset}
  \uses{def:weil-frobenius-orbit-representative,
    def:weil-frobenius-orbit-length}
  \lean{Weil.FiniteOrbitFrobeniusData.Orbit.pointFinset}

  For an orbit $O$, put $x=\operatorname{rep}(O)$ and $d=d_F(x)$, and define
  \[
    \mathcal P(O):=
      \bigl(\operatorname{Finset.range}d\bigr)
        .\operatorname{image}\bigl(j\longmapsto F^{[j]}(x)\bigr).
  \]
\end{definition}

\begin{lemma}[An integer iterate has a bounded natural representative]
  \label{lemma:weil-frobenius-zpow-bounded-representative}
  \uses{def:weil-frobenius-orbit-length,
    lemma:weil-frobenius-orbit-length-positive}
  \lean{Weil.FiniteOrbitFrobeniusData.zpow_eq_iterate_val}

  Let $x\in\Omega$, put $d=d_F(x)$, and let $z\in\mathbf Z$.  There is a
  natural number $j<d$ such that
  \[
    (F^z)(x)=F^{[j]}(x).
  \]
\end{lemma}

\begin{proof}
  \uses{def:weil-frobenius-orbit-length,
    lemma:weil-frobenius-orbit-length-positive}
  By Lemma~\ref{lemma:weil-frobenius-orbit-length-positive}, $d>0$, so
  \texttt{ZMod d} has the required \texttt{NeZero d} instance.  Let
  $k:=(z:\texttt{ZMod d})$ and $j:=k.\texttt{val}$.  The exact theorem
  \texttt{ZMod.val\_lt} gives $j<d$, while
  \texttt{ZMod.natCast\_zmod\_val} gives
  $(j:\texttt{ZMod d})=k=(z:\texttt{ZMod d})$.

  Apply the inverse of \texttt{MulAction.orbitZPowersEquiv F x} to this
  equality.  Its exact evaluation theorem
  \texttt{MulAction.orbitZPowersEquiv\_symm\_apply'} identifies the two
  sides, after taking underlying points, with $(F^j)(x)$ and $(F^z)(x)$.
  Finally \texttt{zpow\_natCast} and \texttt{Equiv.Perm.coe\_pow} rewrite
  $(F^j)(x)$ as the natural iterate $F^{[j]}(x)$.
\end{proof}

\begin{lemma}[The explicit orbit finset has the correct members]
  \label{lemma:weil-frobenius-orbit-point-finset-mem}
  \uses{def:weil-frobenius-orbit-point-finset,
    def:weil-frobenius-orbit-type}
  \lean{Weil.FiniteOrbitFrobeniusData.Orbit.mem_pointFinset}

  For $y\in\Omega$ and an orbit $O$,
  \[
    y\in\mathcal P(O)\quad\Longleftrightarrow\quad y\in|O|.
  \]
\end{lemma}

\begin{proof}
  \uses{lemma:weil-frobenius-orbit-length-positive,
    lemma:weil-frobenius-orbit-representative-class,
    lemma:weil-frobenius-zpow-bounded-representative}
  Put $x=\operatorname{rep}(O)$ and $d=d_F(x)$.  A member of
  $\mathcal P(O)$ has the form $F^{[j]}(x)$, so the integer witness
  $(j:\mathbf Z)$ puts it in the
  $\operatorname{Subgroup.zpowers}(F)$-orbit of $x$ and hence in $|O|$.

  Conversely, rewrite $|O|$ as the orbit of $x$ by
  \texttt{MulAction.orbitRel.Quotient.orbit\_eq\_orbit\_out}.  A point in
  this orbit has the form $F^z(x)$ for $z\in\mathbf Z$.  Apply
  Lemma~\ref{lemma:weil-frobenius-zpow-bounded-representative} to obtain
  $j<d$ with $F^z(x)=F^{[j]}(x)$.  Hence $j$ belongs to
  $\operatorname{Finset.range}d$ and $y$ belongs to its image.
\end{proof}

\begin{lemma}[Cardinality of the explicit orbit finset]
  \label{lemma:weil-frobenius-orbit-point-finset-card}
  \uses{def:weil-frobenius-orbit-point-finset,
    def:weil-frobenius-orbit-degree}
  \lean{Weil.FiniteOrbitFrobeniusData.Orbit.card_pointFinset}

  For every orbit $O$,
  \[
    \#\mathcal P(O)=\deg_F(O).
  \]
\end{lemma}

\begin{proof}
  \uses{lemma:weil-frobenius-orbit-point-finset-mem,
    lemma:weil-frobenius-orbit-finite}
  Lemma~\ref{lemma:weil-frobenius-orbit-point-finset-mem} identifies the
  coercion of $\mathcal P(O)$ to a set with $|O|$.  For a finite set, the
  cardinal of any finset having exactly that membership predicate is its
  \texttt{Set.ncard}.  The latter is $\deg_F(O)$ by definition.
\end{proof}

\begin{lemma}[Distinct orbits have disjoint point finsets]
  \label{lemma:weil-frobenius-orbit-point-finsets-disjoint}
  \uses{def:weil-frobenius-orbit-point-finset}
  \lean{Weil.FiniteOrbitFrobeniusData.Orbit.pointFinset_disjoint}

  If $O\ne O'$, then
  $\operatorname{Disjoint}(\mathcal P(O),\mathcal P(O'))$.
\end{lemma}

\begin{proof}
  \uses{lemma:weil-frobenius-orbit-point-finset-mem}
  Suppose $y$ belonged to both finsets.  Then $y\in|O|$ and $y\in|O'|$.
  The exact equivalence
  \nolinkurl{MulAction.orbitRel.Quotient.mem_orbit} turns these memberships
  into $[y]_F=O$ and $[y]_F=O'$.  Transitivity gives $O=O'$, a
  contradiction.
\end{proof}

\begin{definition}[Orbit class of a point]
  \label{def:weil-frobenius-orbit-class}
  \uses{def:weil-frobenius-orbit-type}
  \lean{Weil.FiniteOrbitFrobeniusData.orbitClass}

  Define $\operatorname{cl}_F(x):=[x]_F$.
\end{definition}

\begin{definition}[Finset of orbits of bounded degree]
  \label{def:weil-bounded-frobenius-orbit-finset}
  \uses{def:weil-finite-orbit-frobenius-datum,
    def:weil-frobenius-fixed-point-finset,
    def:weil-frobenius-orbit-class,
    def:weil-frobenius-orbit-degree}
  \lean{Weil.FiniteOrbitFrobeniusData.boundedOrbitFinset}

  For every $r\in\mathbf N$, define
  \[
    \mathcal O_{\leq r}:=
      \left(
        \operatorname{Fix}_D(r!)
          .\operatorname{image}(\operatorname{cl}_F)
      \right)
      .\operatorname{filter}\bigl(O\longmapsto\deg_F(O)\leq r\bigr).
  \]
\end{definition}

\begin{lemma}[Membership criterion for bounded orbit finsets]
  \label{lemma:weil-bounded-frobenius-orbit-finset-mem}
  \uses{def:weil-bounded-frobenius-orbit-finset}
  \lean{Weil.FiniteOrbitFrobeniusData.mem_boundedOrbitFinset}

  If $r\geq1$, then
  \[
    O\in\mathcal O_{\leq r}\quad\Longleftrightarrow\quad
    \deg_F(O)\leq r.
  \]
\end{lemma}

\begin{proof}
  \uses{lemma:weil-frobenius-orbit-degree-positive,
    lemma:weil-frobenius-orbit-representative-class,
    lemma:weil-frobenius-fixed-iff-orbit-degree-divides,
    lemma:weil-frobenius-fixed-point-finset-mem}
  Forward implication is the filter condition.  Conversely assume
  $1\leq\deg_F(O)\leq r$.  The exact arithmetic theorem
  \texttt{Nat.dvd\_factorial} gives $\deg_F(O)\mid r!$.  Apply
  Lemma~\ref{lemma:weil-frobenius-fixed-iff-orbit-degree-divides} to
  $x=\operatorname{rep}(O)$ and $n=r!$, and use
  $[\operatorname{rep}(O)]_F=O$.  This gives
  $F^{[r!]}(\operatorname{rep}(O))=\operatorname{rep}(O)$.
  Since \texttt{Nat.factorial\_pos r} gives $0<r!$,
  Lemma~\ref{lemma:weil-frobenius-fixed-point-finset-mem} converts that
  equation into
  $\operatorname{rep}(O)\in\operatorname{Fix}_D(r!)$.  Its image under
  $\operatorname{cl}_F$ is $O$,
  so $O$ belongs to the image finset, and the assumed inequality puts it in
  the filter.
\end{proof}

\begin{definition}[Finset of orbits of one degree]
  \label{def:weil-degree-frobenius-orbit-finset}
  \uses{def:weil-bounded-frobenius-orbit-finset,
    def:weil-frobenius-orbit-degree}
  \lean{Weil.FiniteOrbitFrobeniusData.degreeOrbitFinset}

  For every $m\in\mathbf N$, define
  \[
    \mathcal O_m:=
      \mathcal O_{\leq m}
        .\operatorname{filter}\bigl(O\longmapsto\deg_F(O)=m\bigr).
  \]
\end{definition}

\begin{lemma}[Membership criterion for one-degree orbit finsets]
  \label{lemma:weil-degree-frobenius-orbit-finset-mem}
  \uses{def:weil-degree-frobenius-orbit-finset}
  \lean{Weil.FiniteOrbitFrobeniusData.mem_degreeOrbitFinset}

  If $m\geq1$, then
  \[
    O\in\mathcal O_m\quad\Longleftrightarrow\quad\deg_F(O)=m.
  \]
\end{lemma}

\begin{proof}
  \uses{lemma:weil-bounded-frobenius-orbit-finset-mem}
  Forward implication is the second filter condition.  Conversely,
  $\deg_F(O)=m$ implies $\deg_F(O)\leq m$, so
  Lemma~\ref{lemma:weil-bounded-frobenius-orbit-finset-mem} puts $O$ in
  $\mathcal O_{\leq m}$; the equality then puts it in the filter.
\end{proof}

\begin{definition}[Number of Frobenius orbits of one degree]
  \label{def:weil-frobenius-orbit-count}
  \uses{def:weil-degree-frobenius-orbit-finset}
  \lean{Weil.FiniteOrbitFrobeniusData.orbitCount}

  For every $m\in\mathbf N$, define
  \[
    a_m(D):=\#\mathcal O_m.
  \]
\end{definition}

\begin{definition}[Relevant orbits for a Frobenius power]
  \label{def:weil-relevant-frobenius-orbit-finset}
  \uses{def:weil-bounded-frobenius-orbit-finset,
    def:weil-frobenius-orbit-degree}
  \lean{Weil.FiniteOrbitFrobeniusData.relevantOrbitFinset}

  For every $n\in\mathbf N$, define
  \[
    \mathcal O(n):=
      \mathcal O_{\leq n}
        .\operatorname{filter}\bigl(O\longmapsto\deg_F(O)\mid n\bigr).
  \]
\end{definition}

\begin{lemma}[Membership criterion for relevant orbit finsets]
  \label{lemma:weil-relevant-frobenius-orbit-finset-mem}
  \uses{def:weil-relevant-frobenius-orbit-finset}
  \lean{Weil.FiniteOrbitFrobeniusData.mem_relevantOrbitFinset}

  If $n\geq1$, then
  \[
    O\in\mathcal O(n)\quad\Longleftrightarrow\quad\deg_F(O)\mid n.
  \]
\end{lemma}

\begin{proof}
  \uses{lemma:weil-bounded-frobenius-orbit-finset-mem,
    lemma:weil-frobenius-orbit-degree-positive}
  Forward implication is the filter condition.  Conversely, if
  $\deg_F(O)\mid n$, positivity of the degree and $n\geq1$ imply
  $\deg_F(O)\leq n$.  The bounded-orbit membership lemma puts $O$ in
  $\mathcal O_{\leq n}$, and divisibility puts it in the filter.
\end{proof}

\begin{definition}[Orbitwise union of fixed points]
  \label{def:weil-fixed-point-orbit-union}
  \uses{def:weil-relevant-frobenius-orbit-finset,
    def:weil-frobenius-orbit-point-finset}
  \lean{Weil.FiniteOrbitFrobeniusData.fixedPointOrbitUnion}

  For every $n\in\mathbf N$, define the finset
  \[
    U_n:=\mathcal O(n).\operatorname{biUnion}
       \bigl(O\longmapsto\mathcal P(O)\bigr).
  \]
\end{definition}

\begin{lemma}[The orbitwise union is the fixed-point finset]
  \label{lemma:weil-fixed-point-orbit-union-eq}
  \uses{def:weil-finite-orbit-frobenius-datum,
    def:weil-frobenius-fixed-point-finset,
    def:weil-fixed-point-orbit-union}
  \lean{Weil.FiniteOrbitFrobeniusData.fixedPointOrbitUnion_eq}

  For every $n\geq1$,
  \[
    U_n=\operatorname{Fix}_D(n).
  \]
\end{lemma}

\begin{proof}
  \uses{lemma:weil-relevant-frobenius-orbit-finset-mem,
    lemma:weil-frobenius-orbit-point-finset-mem,
    lemma:weil-frobenius-fixed-iff-orbit-degree-divides,
    lemma:weil-frobenius-fixed-point-finset-mem}
  Use finset extensionality and fix $x\in\Omega$.  Membership in $U_n$ means
  that there is an $O\in\mathcal O(n)$ with $x\in\mathcal P(O)$.  The point
  membership lemma forces $O=[x]_F$, and relevant-orbit membership is then
  $\deg_F([x]_F)\mid n$.  By the fixed-point criterion this is equivalent to
  $F^{[n]}(x)=x$, which is equivalent to
  $x\in\operatorname{Fix}_D(n)$ by
  Lemma~\ref{lemma:weil-frobenius-fixed-point-finset-mem}.
\end{proof}

\begin{definition}[Number of fixed points]
  \label{def:weil-frobenius-fixed-point-count}
  \uses{def:weil-frobenius-fixed-point-finset}
  \lean{Weil.FiniteOrbitFrobeniusData.fixedPointCount}

  For every $n\in\mathbf N$, define
  \[
    N_n(D):=\#\operatorname{Fix}_D(n).
  \]
\end{definition}

\begin{lemma}[Fixed points counted by relevant orbits]
  \label{lemma:weil-fixed-point-count-relevant-orbits}
  \uses{def:weil-frobenius-fixed-point-count,
    def:weil-relevant-frobenius-orbit-finset,
    def:weil-frobenius-orbit-degree}
  \lean{Weil.FiniteOrbitFrobeniusData.fixedPointCount_eq_sum_orbits}

  For $n\geq1$,
  \[
    N_n(D)=\sum_{O\in\mathcal O(n)}\deg_F(O).
  \]
\end{lemma}

\begin{proof}
  \uses{lemma:weil-fixed-point-orbit-union-eq,
    lemma:weil-frobenius-orbit-point-finsets-disjoint,
    lemma:weil-frobenius-orbit-point-finset-card}
  Replace $\operatorname{Fix}_D(n)$ by $U_n$.  The family
  $O\mapsto\mathcal P(O)$ is pairwise disjoint on $\mathcal O(n)$ by
  Lemma~\ref{lemma:weil-frobenius-orbit-point-finsets-disjoint}.  Apply the
  exact theorem \texttt{Finset.card\_biUnion}; its right side is
  $\sum_{O\in\mathcal O(n)}\#\mathcal P(O)$.  Replace each cardinal by
  $\deg_F(O)$ using
  Lemma~\ref{lemma:weil-frobenius-orbit-point-finset-card}.
\end{proof}

\begin{lemma}[Relevant orbits as a sigma type over divisors]
  \label{lemma:weil-relevant-orbit-degree-sigma-equiv}
  \uses{def:weil-relevant-frobenius-orbit-finset,
    def:weil-degree-frobenius-orbit-finset}
  \lean{Weil.FiniteOrbitFrobeniusData.relevantOrbitDegreeSigmaEquiv}

  If $0<n$, there is an equivalence
  \[
    \{O\mid O\in\mathcal O(n)\}
      \simeq
    \sum_{m:\{m\mid m\in n.\texttt{divisors}\}}
      \{O\mid O\in\mathcal O_m\}.
  \]
  It sends $O$ to the pair $(\deg_F(O),O)$.
\end{lemma}

\begin{proof}
  \uses{lemma:weil-relevant-frobenius-orbit-finset-mem,
    lemma:weil-degree-frobenius-orbit-finset-mem,
    lemma:weil-frobenius-orbit-degree-positive}
  Relevant-orbit membership gives $\deg_F(O)\mid n$; together with positivity
  and $0<n$, \texttt{Nat.mem\_divisors} makes $\deg_F(O)$ an element of
  $n.\texttt{divisors}$.  The one-degree membership criterion puts $O$ in
  the corresponding fiber.

  Conversely, a sigma element consists of
  $m\in n.\texttt{divisors}$ and $O\in\mathcal O_m$.
  By \texttt{Nat.mem\_divisors}, $m\mid n$; since $0<n$,
  \texttt{Nat.pos\_of\_dvd\_of\_pos} gives $0<m$.  Only now apply
  Lemma~\ref{lemma:weil-degree-frobenius-orbit-finset-mem} with this
  positivity proof to obtain $\deg_F(O)=m$.  Thus
  $\deg_F(O)\mid n$, and
  Lemma~\ref{lemma:weil-relevant-frobenius-orbit-finset-mem} places $O$ in
  $\mathcal O(n)$.  These maps are inverse because the first coordinate is
  forced by $\deg_F(O)=m$; \texttt{Subtype.ext} and \texttt{Sigma.ext}
  remove the proof fields.
\end{proof}

\begin{lemma}[Regrouping relevant orbits by degree]
  \label{lemma:weil-relevant-orbits-regroup-by-degree}
  \uses{def:weil-relevant-frobenius-orbit-finset,
    def:weil-frobenius-orbit-count}
  \lean{Weil.FiniteOrbitFrobeniusData.sum_relevant_eq_sum_divisors}

  For $n\geq1$,
  \[
    \sum_{O\in\mathcal O(n)}\deg_F(O)
      =\sum_{m\mid n}m\,a_m(D).
  \]
  Here $\sum_{m\mid n}$ means the finite sum
  $\sum_{m\in n.\texttt{divisors}}$; since $0<n$, these are precisely the
  positive divisors of $n$.
\end{lemma}

\begin{proof}
  \uses{lemma:weil-relevant-orbit-degree-sigma-equiv,
    def:weil-frobenius-orbit-count,
    lemma:weil-degree-frobenius-orbit-finset-mem}
  Regard each finset sum as a sum over its membership subtype.  Apply
  \texttt{Fintype.sum\_equiv} to the equivalence of
  Lemma~\ref{lemma:weil-relevant-orbit-degree-sigma-equiv}; the summand on a
  sigma element $(m,O)$ is $\deg_F(O)$.  Rewrite the sigma-type sum over
  \texttt{Finset.univ}; then use the exact identity
  \texttt{Finset.univ\_sigma\_univ} in reverse and apply
  \texttt{Finset.sum\_sigma}.  Finally, \texttt{Finset.sum\_attach} for the
  outer divisor finset and each orbit finset rewrites the result as
  \[
    \sum_{m\in n.\texttt{divisors}}
      \sum_{O\in\mathcal O_m}\deg_F(O).
  \]
  In the $m$-fiber,
  Lemma~\ref{lemma:weil-degree-frobenius-orbit-finset-mem} gives
  $\deg_F(O)=m$, so
  \texttt{Finset.sum\_const} evaluates the inner sum as
  $m\,\#\mathcal O_m=m\,a_m(D)$.  Substitution in the outer finite sum is
  the displayed equality.
\end{proof}

\begin{proposition}[Fixed points counted orbit by orbit]
  \label{prop:weil-fixed-points-orbit-count}
  \uses{def:weil-frobenius-fixed-point-count,
    def:weil-frobenius-orbit-count}
  \lean{Weil.FiniteOrbitFrobeniusData.fixedPointCount_eq_divisorSum}

  For every $n\geq1$,
  \[
    N_n(D)=\sum_{m\mid n}m\,a_m(D).
  \]
\end{proposition}

\begin{proof}
  \uses{lemma:weil-fixed-point-count-relevant-orbits,
    lemma:weil-relevant-orbits-regroup-by-degree}
  Substitute the right side of
  Lemma~\ref{lemma:weil-fixed-point-count-relevant-orbits} into
  Lemma~\ref{lemma:weil-relevant-orbits-regroup-by-degree}.
\end{proof}

\begin{definition}[Möbius numerator of an orbit count]
  \label{def:weil-orbit-count-mobius-numerator}
  \uses{def:weil-frobenius-fixed-point-count}
  \lean{Weil.FiniteOrbitFrobeniusData.mobiusNumerator}

  For every $m\in\mathbf N$, define the integer
  \[
    M_m(D):=
      \sum_{r\in m.\texttt{divisors}}
        \mu(r)\,(N_{m/r}(D):\mathbf Z),
  \]
  where $\mu=\texttt{ArithmeticFunction.moebius}$.
\end{definition}

\begin{lemma}[Reindexing the divisors antidiagonal]
  \label{lemma:weil-divisors-antidiagonal-reindex}
  \uses{}
  \lean{Weil.sum_divisorsAntidiagonal_moebius}

  If $0<n$ and $g:\mathbf N\to\mathbf Z$, then
  \[
    \sum_{(r,s)\in n.\texttt{divisorsAntidiagonal}}
      \mu(r)g(s)
      =
    \sum_{r\in n.\texttt{divisors}}
      \mu(r)g(n/r).
  \]
\end{lemma}

\begin{proof}
  \uses{}
  Instantiate the exact Mathlib theorem
  \texttt{Nat.sum\_divisorsAntidiagonal} with the two-variable function
  \[
    f(r,s):=(\mu(r):\mathbf Z)\,g(s).
  \]
  Its left and right sides are definitionally the two displayed finite sums.
  The hypothesis $0<n$ is not needed by that library theorem, but is retained
  here because all later Möbius applications are at positive indices.
\end{proof}

\begin{lemma}[Möbius numerator equals degree times orbit count]
  \label{lemma:weil-orbit-count-mobius-numerator-eq}
  \uses{def:weil-orbit-count-mobius-numerator,
    def:weil-frobenius-orbit-count}
  \lean{Weil.FiniteOrbitFrobeniusData.mobiusNumerator_eq}

  For $m\geq1$,
  \[
    M_m(D)=(m:\mathbf Z)\,(a_m(D):\mathbf Z)
  \]
  as an equality of integers.
\end{lemma}

\begin{proof}
  \uses{prop:weil-fixed-points-orbit-count,
    lemma:weil-divisors-antidiagonal-reindex}
  Define total functions $f,g:\mathbf N\to\mathbf Z$ by
  \[
    f(i):=(i:\mathbf Z)(a_i(D):\mathbf Z),
    \qquad g(n):=(N_n(D):\mathbf Z).
  \]
  For every $0<n$, casting
  the reverse of Proposition~\ref{prop:weil-fixed-points-orbit-count} by
  \[
    \texttt{congrArg}\ (\lambda z:\mathbf N,\,(z:\mathbf Z))
  \]
  gives an equality in $\mathbf Z$.  Rewriting it with
  \texttt{Nat.cast\_sum} and \texttt{Nat.cast\_mul} gives exactly
  $\sum_{i\in n.\texttt{divisors}}f(i)=g(n)$.
  Apply the forward direction of
  \texttt{ArithmeticFunction.sum\_eq\_iff\_sum\_mul\_moebius\_eq} with
  coefficient ring $\mathbf Z$ and these total functions.  Its conclusion is
  \[
    \sum_{(r,s)\in n.\texttt{divisorsAntidiagonal}}
       \mu(r)g(s)=f(n).
  \]
  Apply Lemma~\ref{lemma:weil-divisors-antidiagonal-reindex}.  At $n=m$ its
  right side is exactly $M_m(D)$, while $f(m)$ is
  $(m:\mathbf Z)(a_m(D):\mathbf Z)$.
\end{proof}

\begin{corollary}[Möbius recovery of orbit counts]
  \label{cor:weil-orbit-count-mobius-recovery}
  \uses{def:weil-frobenius-orbit-count,
    def:weil-frobenius-fixed-point-count}
  \lean{Weil.FiniteOrbitFrobeniusData.orbitCount_eq_mobius}

  For $m\geq1$,
  \[
    (a_m(D):\mathbf Q)=
      \frac{1}{(m:\mathbf Q)}
      \sum_{r\in m.\texttt{divisors}}
        (\mu(r):\mathbf Q)(N_{m/r}(D):\mathbf Q)
  \]
  in $\mathbf Q$.
\end{corollary}

\begin{proof}
  \uses{def:weil-orbit-count-mobius-numerator,
    lemma:weil-orbit-count-mobius-numerator-eq}
  Apply
  $\texttt{congrArg}\ (\lambda z:\mathbf Z,\,(z:\mathbf Q))$ to the integer
  equality $M_m(D)=(m:\mathbf Z)(a_m(D):\mathbf Z)$.  Rewriting by
  \texttt{Int.cast\_mul} identifies its right side with
  $(m:\mathbf Q)(a_m(D):\mathbf Q)$.  The hypothesis $m\geq1$ gives
  $(m:\mathbf Q)\ne0$, so multiply by its inverse and cancel this factor.
  Finally expand $M_m(D)$; \texttt{Int.cast\_sum} and
  \texttt{Int.cast\_mul} turn its cast into
  \[
    \sum_{r\in m.\texttt{divisors}}
      (\mu(r):\mathbf Q)(N_{m/r}(D):\mathbf Q),
  \]
  which is the displayed formula.
\end{proof}

\begin{corollary}[Divisibility and nonnegativity of Möbius recovery]
  \label{cor:weil-orbit-count-mobius-integrality}
  \uses{def:weil-frobenius-fixed-point-count}
  \lean{Weil.FiniteOrbitFrobeniusData.mobius_divisible_nonnegative}

  For $m\geq1$, the integer
  $M_m(D)=\sum_{r\in m.\texttt{divisors}}
    \mu(r)(N_{m/r}(D):\mathbf Z)$ satisfies
  \[
    (m:\mathbf Z)\mid M_m(D),
    \qquad
    M_m(D)/(m:\mathbf Z)=(a_m(D):\mathbf Z)\geq0.
  \]
\end{corollary}

\begin{proof}
  \uses{def:weil-orbit-count-mobius-numerator,
    lemma:weil-orbit-count-mobius-numerator-eq}
  Lemma~\ref{lemma:weil-orbit-count-mobius-numerator-eq} writes the numerator as
  $(m:\mathbf Z)(a_m(D):\mathbf Z)$, proving divisibility.  Since
  $m>0$, one has $(m:\mathbf Z)\ne0$.  Instantiate the exact theorem
  \texttt{Int.mul\_ediv\_cancel\_left} with
  $a=(m:\mathbf Z)$ and $b=(a_m(D):\mathbf Z)$ to obtain
  \[
    ((m:\mathbf Z)(a_m(D):\mathbf Z))/(m:\mathbf Z)
      =(a_m(D):\mathbf Z).
  \]
  Substitution gives the quotient equality, and
  \texttt{Int.ofNat\_nonneg} gives its nonnegativity.
\end{proof}

\end{document}
